% ============================================================
%  INTRODUCTION GÉNÉRALE
% ============================================================
\chapter*{Introduction générale}
\addcontentsline{toc}{chapter}{Introduction générale}
\markboth{Introduction générale}{Introduction générale}

\hspace{1cm}La performance industrielle d'un groupe minier et chimique repose
autant sur ses procédés que sur sa capacité à disposer, au bon moment et au bon
endroit, des matières et des pièces nécessaires à son exploitation. La
\gls{sc} constitue à ce titre un levier de compétitivité majeur : une pièce
manquante peut immobiliser une ligne de production entière, tandis qu'un stock
excédentaire immobilise du capital et occupe de l'espace sans créer de valeur.

\vspace{0.5cm}
Le présent travail s'inscrit dans le cadre du \gls{pfa} pour la formation
d'Ingénieur d'État en \textbf{\gls{iacs}} à l'École Nationale des Sciences
Appliquées de Béni Mellal, réalisé au sein du groupe \gls{ocp}.

\vspace{0.5cm}
La problématique abordée est celle de l'arbitrage permanent entre
\textbf{disponibilité} et \textbf{coût} : comment dimensionner les stocks et les
règles de réapprovisionnement de manière à répondre à la demande de
l'exploitation, tout en maîtrisant les coûts, la qualité et les volumes ? Cet
arbitrage est rendu difficile par deux sources d'incertitude : la variabilité de
la consommation et celle des délais fournisseurs.

\vspace{0.5cm}
Pour y répondre, la démarche adoptée consiste à exploiter l'historique des
consommations et des stocks afin de segmenter le portefeuille de références,
caractériser statistiquement la demande, puis calculer et optimiser les
paramètres de gestion (stock de sécurité, point de commande, quantité de
commande). Les résultats sont restitués sous la forme d'un tableau de bord
d'aide à la décision.

\vspace{0.5cm}
Ce rapport est organisé en trois chapitres principaux :

\begin{itemize}
  \item \textbf{Chapitre I -- Contexte général et problématique :} présentation
    du groupe \gls{ocp} et de son organisation logistique, description de la
    situation actuelle de la gestion des stocks, formulation de la
    problématique et du cahier des charges.
  \item \textbf{Chapitre II -- Cadre théorique et état de l'art :} rappel des
    modèles classiques de gestion des stocks, des méthodes de classification
    du portefeuille de références et des techniques de prévision de la demande.
  \item \textbf{Chapitre III -- Conception, réalisation et résultats :}
    construction de la chaîne de traitement des données, mise en \oe uvre des
    modèles, développement du tableau de bord et analyse des gains obtenus.
\end{itemize}
